\ifdefined\gkver\else\def\gkver{student}\fi
\documentclass[\gkver]{gaokaozhenti}
\begin{document}

\chapter{2008年重庆理综}
%% paperType: 理综物理部分

\begin{enumerate}
\item
%% number: 14
%% paperName: 2008年重庆理综第 14 题 6 分
%% typeId: 1
%% type: 单选题
%% chapter: 原子和原子核
%% point: 天然放射性
%% method: 无
%% score: 6
%% degree: 850
%% duplicateId: 0
%% body:
放射性同位素钍 232 经 $\alpha$、$\beta$ 衰变会生成氡，其衰变方程为 $\ce{^232_90 Th -> ^220_86 Rn}+x\alpha+y\beta$，其中\xzanswer{D}

\fourchoices[answer=D]
{$x=1$，$y=3$}
{$x=2$，$y=3$}
{$x=3$，$y=1$}
{$x=3$，$y=2$}

%% answer: D
%% memo:
\memoanswer{本题考查放射性元素衰变的有关知识。由衰变方程 $\ce{^232_90 Th -> ^220_86 Rn}+x\ce{^4_2 He}+y\ce{^0_{-1} e}$，由质量数守恒和电荷数守恒得：$232=220+4x$，$90=86+2x-y$，可解得：$x=3$、$y=2$。}

\item
%% number: 15
%% paperName: 2008年重庆理综第 15 题 6 分
%% typeId: 1
%% type: 单选题
%% chapter: 电路
%% point: 串并联电路
%% method: 无
%% score: 6
%% degree: 550
%% duplicateId: 0
%% body:
某同学设计了一个转向灯电路，其中 L 为指示灯，$L_{1}$、$L_{2}$ 分别为左、右转向灯，S 为单刀双掷开关，$E$ 为电源，当 S 置于位置 1 时，以下判断正确的是\xzanswer{A}

\onepicture[w=4.5cm]{figs/15.png}

\fourchoices[answer=A]
{L 的功率小于额定功率}
{$L_{1}$ 亮，其功率等于额定功率}
{$L_{2}$ 亮，其功率等于额定功率}
{含 L 支路总功率较另一支路大}

%% answer: A
%% memo:
\memoanswer{本题考查电路分析的有关知识。由电路结构可知，当 S 置于位置 1 时，L 与 $L_{2}$ 串联后再与 $L_{1}$ 并联，由灯泡的额定电压和额定功率可知，$L_{1}$ 和 $L_{2}$ 的电阻相等。L 与 $L_{2}$ 串联后的总电阻大于 $L_{1}$ 的电阻，由于电源电动势为 $6\UV$，本身有电阻，所以 $L_{1}$ 两端电压和 L 与 $L_{2}$ 的总电压相等，且都小于 $6\UV$，所以三只灯都没有正常发光，三只灯的实际功率都小于额定功率。含 L 的支路的总电阻大于 $L_{1}$ 支路的电阻，由于两条支路的电压相等，所以含 L 的支路的总功率小于另一支路的功率。}

\item
%% number: 16
%% paperName: 2008年重庆理综第 16 题 6 分
%% typeId: 1
%% type: 单选题
%% chapter: 气体性质
%% point: 热力学第一定律
%% method: 无
%% score: 6
%% degree: 650
%% duplicateId: 0
%% body:
地面附近有一正在上升的空气团，它与外界的热交换忽略不计，已知大气压强随高度增加而降低，则该气团在此上升过程中（不计气团内分子间的势能）\xzanswer{C}

\fourchoices[answer=C]
{体积减小，温度降低}
{体积减小，温度不变}
{体积增大，温度降低}
{体积增大，温度不变}

%% answer: C
%% memo:
\memoanswer{本题考查气体的有关知识。随着空气团的上升，大气压强也随着减小，那么空气团的体积会增大，空气团对外做功，其内能会减小，因为不计分子势能，所以内能由其温度决定，则其温度会降低。所以空气团的体积增大、温度降低、压强减小。}

\item
%% number: 17
%% paperName: 2008年重庆理综第 17 题 6 分
%% typeId: 1
%% type: 单选题
%% chapter: 功和能
%% point: 能量守恒定律
%% method: 无
%% score: 6
%% degree: 650
%% duplicateId: 0
%% body:
下列与能量有关的说法中正确的是\xzanswer{B}

\fourchoices[answer=B]
{卫星绕地球圆周运动的半径越大，动能越大}
{从同种金属逸出的光电子的最大初动能随照射光的波长增大而减小}
{做平抛运动的物体在任意相等时间内动能的增量相同}
{在静电场中，电场线越密的地方正电荷的电势能越大}

%% answer: B
%% memo:
\memoanswer{本题考查能量有关的问题。卫星绕地球做圆周运动中，$v=\sqrt{\frac{GM}{r}}$，半径越大，其速度越小，其动能也就越小；根据光电效应方程 $h\frac{c}{\lambda}=W+E_{k}$，对于同一种金属而言，$W$ 是一定的，所以入射光的波长减小，其最大初动能增大；做平抛运动的物体，其动能的变化量为 $\Delta E_{k}=\frac{1}{2}m(v_{y2}^{2}-v_{y1}^{2})=\frac{1}{2}m\{(v_{y1}+gt)^{2}-v_{y1}^{2}\}=mv_{y1}gt+\frac{1}{2}mg^{2}t^{2}$，所以任意相等时间内动能的增量不相等；在静电场中，电场线越密的地方，其电场强度越大，但其电势可为正，也可为负，所以正电荷在电场线越密的地方，其电势能并不一定越大。}

\item
%% number: 18
%% paperName: 2008年重庆理综第 18 题 6 分
%% typeId: 1
%% type: 单选题
%% chapter: 电磁感应
%% point: 楞次定律推广
%% method: 无
%% score: 6
%% degree: 650
%% duplicateId: 0
%% body:
如图，粗糙水平桌面上有一质量为 $m$ 的铜质矩形线圈，当一竖直放置的条形磁铁从线圈中线 AB 正上方等高快速经过时，若线圈始终不动，则关于线圈受到的支持力 $F_{N}$ 及在水平方向运动趋势的正确判断是\xzanswer{D}

\onepicture[w=4.0cm]{figs/18.png}

\fourchoices[answer=D]
{$F_{N}$ 先小于 $mg$ 后大于 $mg$，运动趋势向左}
{$F_{N}$ 先大于 $mg$ 后小于 $mg$，运动趋势向左}
{$F_{N}$ 先小于 $mg$ 后大于 $mg$，运动趋势向右}
{$F_{N}$ 先大于 $mg$ 后小于 $mg$，运动趋势向右}

%% answer: D
%% memo:
\memoanswer{本题考查电磁感应有关的知识。条形磁铁从线圈正上方等高快速经过时，通过线圈的磁通量先增加后又减小。当通过线圈磁通量增加时，为阻碍其增加，在竖直方向上线圈有向下运动的趋势，所以线圈受到的支持力大于其重力，在水平方向上有向右运动的趋势；当通过线圈的磁通量减小时，为阻碍其减小，在竖直方向上线圈有向上运动的趋势，所以线圈受到的支持力小于其重力，在水平方向上有向右运动的趋势。综上所述，线圈所受到的支持力先大于重力后小于重力，运动趋势总是向右。}

\item
%% number: 19
%% paperName: 2008年重庆理综第 19 题 6 分
%% typeId: 1
%% type: 单选题
%% chapter: 光学
%% point: 光的折射
%% method: 无
%% score: 6
%% degree: 650
%% duplicateId: 0
%% body:
图为一个四分之一圆柱体棱镜的截面图，图中 E、F、G、H 将半径 OM 分成五等份，虚线 $EE_{1}$、$FF_{1}$、$GG_{1}$、$HH_{1}$ 平行于半径 ON，ON 边可吸收到达其上的所有光线。已知该棱镜的折射率 $n=\frac{5}{3}$，若平行光束垂直入射并覆盖 OM，则光线\xzanswer{B}

\onepicture[w=3.5cm]{figs/19.png}

\fourchoices[answer=B]
{不能从圆弧 $NF_{1}$ 射出}
{不能从圆弧 $NG_{1}$ 射出}
{不能从圆弧 $G_{1}H_{1}$ 射出}
{不能从圆弧 $H_{1}M$ 射出}

%% answer: B
%% memo:
\memoanswer{本题考查光的折射有关的知识。由该棱镜的折射率为 $n=\frac{5}{3}$ 可知其临界角 $C$ 满足：$\sin C=\frac{1}{n}=\frac{3}{5}$，可求出 $GG_{1}$ 右边的入射光线没有发生全反射，其左边的光线全部发生全反射。所以光线只能从圆弧 $NG_{1}$ 射出。}

\item
%% number: 20
%% paperName: 2008年重庆理综第 20 题 6 分
%% typeId: 1
%% type: 单选题
%% chapter: 机械波
%% point: 波长频率波速
%% method: 无
%% score: 6
%% degree: 650
%% duplicateId: 0
%% body:
某地区地震波中的横波和纵波传播速率分别约为 $4\Ukms$ 和 $9\Ukms$，一种简易地震仪由一竖直弹簧振子 P 和一水平弹簧振子 H 组成，如图。在一次地震中震源在地震仪下方，观察到两振子相差 $5\Us$ 开始振动，则\xzanswer{A}

\onepicture[w=3.5cm]{figs/20.png}

\fourchoices[answer=A]
{P 先开始振动，震源距地震仪约 $36\Ukm$}
{P 先开始振动，震源距地震仪约 $25\Ukm$}
{H 先开始振动，震源距地震仪约 $36\Ukm$}
{H 先开始振动，震源距地震仪约 $25\Ukm$}

%% answer: A
%% memo:
\memoanswer{本题考查地震波有关的知识。由于纵波的传播速度快些，所以纵波先到达地震仪处，所以 P 先开始振动。设地震仪距震源为 $x$，则有 $\frac{x}{4}-\frac{x}{9}=5$，解得 $x=36\Ukm$。}

\item
%% number: 21
%% paperName: 2008年重庆理综第 21 题 6 分
%% typeId: 1
%% type: 单选题
%% chapter: 电路
%% point: 电容
%% method: 无
%% score: 6
%% degree: 450
%% duplicateId: 0
%% body:
图 1 为某同学设计的电容式速度传感器原理图，其中上板为固定极板，下板为待测物体。在两极板间电压恒定的条件下，极板上所带电量 $Q$ 将随待测物体的上下运动而变化，若 $Q$ 随时间 $t$ 的变化关系为 $Q=\frac{b}{t+a}$（$a$、$b$ 为大于零的常数），其图像如图 2 所示，那么图 3、图 4 中反映极板间场强大小 $E$ 和物体速率 $v$ 随时间 $t$ 变化关系可能是\xzanswer{C}

\twopicture[showlabel=false, h=2.5cm]{figs/21a.png}{figs/21b.png}

\twopicture[showlabel=false, h=2.5cm]{figs/21c.png}{figs/21d.png}

\fourchoices[answer=C]
{①和③}
{①和④}
{②和③}
{②和④}

%% answer: C
%% memo:
\memoanswer{本题考查速度传感器的有关知识。由题意可知：$E=\frac{U}{d}=\frac{Q}{Cd}=\frac{Q}{\frac{\varepsilon s}{4k\pi d}\cdot d}=\frac{Q}{\frac{\varepsilon s}{4k\pi}}$，所以 $E$ 的变化规律与 $Q$ 的变化规律相似，故 $E$ 的图象为②；由 $Q=CU=\frac{\varepsilon s}{4k\pi d}U=\frac{b}{t+a}$，得 $d=vt+a$，即物体匀速移动，故速度图象为③。综上所述，C 正确。}

\item
%% number: 22
%% paperName: 2008年重庆理综第 22 题 11 分
%% typeId: 4
%% type: 实验题
%% chapter: 电路
%% point: 电路实验
%% method: 无
%% score: 11
%% degree: 650
%% duplicateId: 0
%% body:
\begin{enumerate}
	\item
	某实验小组拟用图 1 所示装置研究滑块的运动，实验器材有滑块、钩码、纸带、米尺、带滑轮的木板，以及由漏斗和细线组成的单摆等。实验中，滑块在钩码作用下拖动纸带做匀加速直线运动，同时单摆垂直于纸带运动方向摆动，漏斗漏出的有色液体在纸带上留下的痕迹记录了漏斗在不同时刻的位置。
	
	①在图 2 中，从 \tkanswer[1.5cm]{B} 纸带可看出滑块的速度和加速度方向一致；
	
	②用该方法测量滑块加速度的误差主要来源有 \tkanswer[2.2cm]{摆长测量}、\tkanswer[2.2cm]{漏斗重心变化}（写出 2 个即可）。
	
	\twopicture[showlabel=false, h=3.2cm]{figs/22a1.png}{figs/22a2.png}
	
	\item
	某研究性小组设计了如图 1 所示的电路，用来研究稀盐水溶液的电阻率和浓度的关系。图中 $E$ 为直流电源，K 为开关，$K_{1}$ 为单刀双掷开关，V 为电压表，A 为多量程电流表，$R$ 为滑动变阻器，$R_{x}$ 为待测稀盐水溶液液柱。
	
	①实验时，闭合 K 之前应将 $R$ 的滑片置于 \tkanswer[1cm]{D}（填“C”或“D”）端，当用电流表外接法测量 $R_{x}$ 的阻值时，$K_{1}$ 应置于位置 \tkanswer[1cm]{1}（填“1”或“2”）。
	
	②在一定条件下，用电流表内、外接法得到 $R_{x}$ 的电阻率随浓度变化的两条曲线如图 2 所示（不计由于通电导致的化学变化）。实验中 $R_{x}$ 的通电面积为 $20\Ucmq$，长为 $20\Ucm$，用内接法测得 $R_{x}$ 为 $3500\UO$，其电阻率为 \tkanswer[1.2cm]{35} $\UOm$，由图中对应曲线 \tkanswer[1cm]{1}（填“1”或“2”）可得此时溶液浓度约为 \tkanswer[2cm]{0.011～0.014}%（结果保留两位有效数字）。
	
	\twopicture[showlabel=false, h=3.8cm]{figs/22b1.png}{figs/22b2.png}
\end{enumerate}

%% answer:
\jdanswer{
\begin{enumerate}
	\item
	① B；② 摆长测量、漏斗重心变化（或液体痕迹偏粗、阻力变化等）。

	\item
	① D，1；② $35$，1，0.011～0.014 之间都正确。
\end{enumerate}
}

%% memo:
\memoanswer{
（1）本题考查匀加速运动和单摆的运动。要使速度和加速度方向相同，则只有 B，因为 B 中相等时间内纸带运动的距离越来越大。用该方法测量加速度的误差主要来源有：摆长的测量、漏斗重心的变化、液体痕迹偏粗、阻力变化等。

（2）本题考查电阻率测量这个实验。闭合 K 之前，应将 R 的滑片位置置于 D 端，这样可保证电表的安全；用电流表外接法测电阻时，$K_{1}$ 应置于 1 位置。根据电阻定律得 $R=\rho\frac{l}{s}$，得 $\rho=\frac{Rs}{l}=35\UOm$；内接法使得测量值偏大，对应的电阻率也偏大，所以对应的曲线为 1。由测量得的电阻率和图象可得溶液浓度约为 0.011～0.014。
}

\item
%% number: 23
%% paperName: 2008年重庆理综第 23 题 16 分
%% typeId: 6
%% type: 计算题
%% chapter: 力物体的平衡
%% point: 三力平衡
%% method: 无
%% score: 16
%% degree: 650
%% duplicateId: 0
%% body:
滑板运动是一种非常刺激的水上运动。研究表明，在进行滑板运动时，水对滑板的作用力 $F_{N}$ 垂直于板面，大小为 $kv^{2}$，其中 $v$ 为滑板的速率（水可视为静止）。某次运动中，在水平牵引力作用下，当滑板与水面的夹角为 $\theta=37^{\circ}$ 时（如图），滑板做匀速直线运动，相应的 $k=54\Ukg/\Um$，人和滑板的总质量为 $108\Ukg$，试求：（重力加速度 $g$ 取 $10\Umsq$，$\sin 37^{\circ}=0.6$，空气阻力不计）

\begin{enumerate}
	\item
	水平牵引力的大小；
	\item
	滑板的速率；
	\item
	水平牵引力的功率。
\end{enumerate}

\onepicture[w=5.5cm, align=right]{figs/23.png}

%% answer:
\jdanswer{
\begin{enumerate}
	\item
	$F=810\UN$
	
	\item
	$v=5\Ums$
	
	\item
	$P=4050\UW$
\end{enumerate}
}

%% memo:
\memoanswer{
（1）以滑板和运动员为研究对象，其受力如图所示。

\onepicture[w=3.0cm, align=right]{figs/23_an.jpg}

由共点力平衡条件可得 $F_{N}\cos\theta=mg$ ①，$F_{N}\sin\theta=F$ ②，由 ①、② 联立，得 $F=810\UN$。

（2）由 $F_{N}=\frac{mg}{\cos\theta}$，$F_{N}=kv^{2}$，得 $v=\sqrt{\frac{mg}{k\cos\theta}}=5\Ums$。

（3）水平牵引力的功率 $P=Fv=4050\UW$。
}

\item
%% number: 24
%% paperName: 2008年重庆理综第 24 题 9 分
%% typeId: 6
%% type: 计算题
%% chapter: 运动和力
%% point: 动量守恒定律
%% method: 无
%% score: 9
%% degree: 650
%% duplicateId: 0
%% body:
图中有一竖直固定在地面的透气圆筒，筒内有一劲度系数为 $k$ 的轻弹簧，其下端固定，上端连接一质量为 $m$ 的薄滑块，圆筒内壁涂有一种新型智能材料——ER 流体，它对滑块的阻力可调。起初滑块静止，ER 流体对其阻力为零，弹簧的长度为 $L$。现有一质量为 $m$ 的物体从距地面 $2L$ 处自由落下，与滑块碰撞后粘在一起向下运动。为保证滑块做匀减速运动，且下移距离为 $\frac{2mg}{k}$ 时速度减为零，ER 流体对滑块的阻力需随滑块下移而变。试求（忽略空气阻力）：

\begin{enumerate}
	\item
	下落物体与滑块碰撞过程中，系统损失的机械能；
	\item
	滑块向下运动过程中加速度大小；
	\item
	滑块下移距离 $d$ 时 ER 流体对滑块阻力的大小。
\end{enumerate}

\onepicture[w=3.5cm, align=right]{figs/24.png}

%% answer:
\jdanswer{
\begin{enumerate}
	\item
	$\frac{1}{2}mgL$
	
	\item
	$a=\frac{kL}{8m}$
	
	\item
	$F_{ER}=mg+\frac{kL}{4}-kd$
\end{enumerate}
}

%% memo:
\memoanswer{
（1）设物体下落末速度为 $v_{0}$，由机械能守恒定律 $mgL=\frac{1}{2}mv_{0}^{2}$，得 $v_{0}=\sqrt{2gL}$。设碰后共同速度为 $v_{1}$，由动量守恒定律 $2mv_{1}=mv_{0}$，得 $v_{1}=\frac{1}{2}\sqrt{2gL}$。碰撞过程中系统损失的机械能为 $\Delta E=\frac{1}{2}mv_{0}^{2}-\frac{1}{2}\cdot 2mv_{1}^{2}=\frac{1}{2}mgL$。

（2）设加速度大小为 $a$，有 $2as=v_{1}^{2}$，得 $a=\frac{kL}{8m}$。

（3）设弹簧弹力为 $F_{N}$，ER 流体对滑块的阻力为 $F_{ER}$，对整体受力分析如图所示。

\onepicture[w=2.6cm, align=right]{figs/24_an.jpg}

对整体有 $F_{N}+F_{ER}-2mg=2ma$，又 $F_{N}=kx$，$x=d+\frac{mg}{k}$，联立解得 $F_{ER}=mg+\frac{kL}{4}-kd$。
}

\item
%% number: 25
%% paperName: 2008年重庆理综第 25 题 20 分
%% typeId: 6
%% type: 计算题
%% chapter: 磁场
%% point: 带电粒子在磁场中的运动
%% method: 无
%% score: 20
%% degree: 450
%% duplicateId: 0
%% body:
图为一种质谱仪工作原理示意图，在以 O 为圆心、OH 为对称轴、夹角为 $2\alpha$ 的扇形区域内分布着方向垂直于纸面的匀强磁场。对称于 OH 轴的 C、D 分别是离子发射点和收集点，CM 垂直交磁场左边界于 M，且 $OM=d$。现有一正离子束以小发散角（纸面内）从 C 射出，这些离子在 CM 方向的分速度均为 $v_{0}$。若该离子束中比荷为 $\frac{q}{m}$ 的离子都能汇聚到 D。试求：

\begin{enumerate}
	\item
	磁感应强度的大小和方向（提示：可考虑沿 CM 方向运动的离子为研究对象）；
	\item
	离子沿与 CM 成 $\theta$ 角的直线 CN 进入磁场，其轨道半径与在磁场中运动的时间；
	\item
	线段 CM 的长度。
\end{enumerate}

\onepicture[w=6.0cm, align=right]{figs/25.png}

%% answer:
\jdanswer{
\begin{enumerate}
	\item
	$B=\frac{mv_{0}}{qd}$，方向垂直纸面向外
	
	\item
	$R'=\frac{d}{\cos\theta}$，$t=\frac{2(\alpha+\theta)d}{v_{0}}$
	
	\item
	$CM=d\cot\alpha$
\end{enumerate}
}

%% memo:
\memoanswer{
（1）设沿 CM 方向运动的离子在磁场中做圆周运动的轨道半径为 $R$，由洛伦兹力提供向心力 $qv_{0}B=m\frac{v_{0}^{2}}{R}$，又 $R=d$，得 $B=\frac{mv_{0}}{qd}$，方向垂直纸面向外。

（2）设沿 CN 运动的离子速度大小为 $v$，在磁场中的轨道半径为 $R'$，运动时间为 $t$。由 $v\cos\theta=v_{0}$ 得 $v=\frac{v_{0}}{\cos\theta}$，$R'=\frac{mv}{qB}=\frac{d}{\cos\theta}$。离子在磁场中做匀速圆周运动的周期 $T=\frac{2\pi m}{qB}$，则 $t=T\cdot\frac{\alpha+\theta}{\pi}=\frac{2(\alpha+\theta)d}{v_{0}}$。

（3）设圆心为 $A$，过 $A$ 作 $AB\perp OM$，可证明 $NM=BO$。又 $NM=CM\tan\theta$，$BO=AB\cot\alpha=R'\sin\theta\cot\alpha=\frac{d}{\cos\theta}\sin\theta\cot\alpha$，所以 $CM=\frac{BO}{\tan\theta}=d\cot\alpha$。
}

\end{enumerate}

\end{document}
